\documentclass[11pt]{article}

\usepackage[margin=1in]{geometry}
\usepackage{amsmath, amssymb, amsthm}
\usepackage{mathtools}
\usepackage{booktabs}
\usepackage{tabularx}
\usepackage{graphicx}
\graphicspath{{figures/}{../figures/}}
\usepackage[hidelinks]{hyperref}
\usepackage[authoryear,round]{natbib}
\bibpunct{(}{)}{,}{a}{}{,}

\newcommand{\R}{\mathbb{R}}
\newcommand{\E}{\mathbb{E}}
\newcommand{\F}{\mathcal{F}}
\newcommand{\pluspart}[1]{\left[#1\right]_+}

\theoremstyle{plain}
\newtheorem{theorem}{Theorem}
\newtheorem{lemma}{Lemma}
\newtheorem{proposition}{Proposition}
\theoremstyle{definition}
\newtheorem{definition}{Definition}
\newtheorem{assumption}{Assumption}
\theoremstyle{remark}
\newtheorem{remark}{Remark}

\title{Optimal Trading Filters: a Wiener--Hopf Approach}
\author{Matteo Smerlak\thanks{Capital Fund Management, 23 rue de l'Universit\'e, 75007 Paris, France.\ \texttt{matteo.smerlak@cfm.com}}}
\date{}

\begin{document}
\maketitle

\begin{center}
\begin{tabular}{p{0.9\linewidth}}
\textbf{Keywords.} Optimal trading; market impact; propagator model; inventory risk; Wiener--Hopf factorization; linear filters; fractional calculus; adapted stochastic control.\\[2pt]
\textbf{JEL classification.} C61; G11; G12; G14.\\[2pt]
\textbf{MSC classification.} 45E10; 60G35; 91G80; 26A33.
\end{tabular}
\end{center}

\begin{abstract}
A trader holding a predictive signal must weigh the expected gain against the cost of acquiring the position and the risk of holding it. Under the propagator model of market impact, trading-induced price changes involve the entire history of past trades. Existing treatments of the resulting path-wise quadratic program characterize the optimum implicitly, as the solution of an integral equation or an equivalent forward--backward system. In this paper I show how stationary solutions can be computed analytically from the signal and its forecasts using a Wiener-Hopf factorization of the friction operator. For rational frictions --- inventory risk, temporary cost, and exponential resilience --- the trading filter reduces to finitely many exponential moving averages, and we recover the Neuman--Vo\ss{} solution away from horizon boundaries. In the (empirically more realistic) case of scale-free impact kernels, the optimal trading filter decays algebraically rather than exponentially at large frequencies; in the limit of small risk aversion, the target position is a fractional integral of the signal. Unlike its exponential (Obizhaeva-Wang) cousin, this solution never surfs its own impact or generates block trades.
\end{abstract}

% =====================================================================
% BODY
% =====================================================================

\section{Introduction}\label{sec:intro}
\subsection{The quadratic optimal trading problem}
\subsection{The adaptedness constraint}
- With instantaneous impact the solution is trivial: Garleanu-Petersen
- Non-local costs create a difficulty: implement adaptedness constraint
- Literature (starting from OW and GSS and more recently Neuman, Voss and collaborators) has delt with this using dynamic programming and/or integral equations. Typically implicit and hard to extract intuition from.
\subsection{The Wiener-Hopf method}
- history and application in filtering and prediction
- general formulation: factorization of positive operator along a nest, Gohberg-Krein and Avreson
- here applied to stochastic processes, with causality = adaptedness
\subsection{Contribution and outline}
- WH solution in terms of signal forecasts
- formulation as linear filter for stationary signals, two-step WH (friction and signal itself)
- analytical solution for any quadratic friction away from boundaries
- connection with fractional calculus
- estimation of causality gap
- an exact solution on the finite interval (for use as execution policy) with better properties than OW

\section{The Wiener--Hopf solution}\label{sec:interior}

\section{Exact filters}\label{sec:recovery}
\subsection{Exponential kernels: zero, and two EMAs}
\subsection{Power-law kernels and fractional calculus}


\section{Finite horizon solutions}\label{sec:recovery}
- general case treated in the literature (NV, AJN, Forde et al, etc), but not in closed form execpt in the absence of signal (OW and GSS)
- we give closed forms for pure impact (exponential or power-law) with signal

\subsection{Gohberg-Krein generalization of WH}
- explain GK 
- bound: WH solution + small away from boundaries
\subsection{Pure transient impact with signal}

- exponential (OW with signal)
- power-law (GSS with signal)

\section{Conclusion}

% \appendix

% =====================================================================
\bibliographystyle{plainnat}
\bibliography{optimal-trading-filters}

\end{document}
